\section{Wprowadzenie}
\label{sec:Wprowadzenie}

\subsection{Geneza}

Pomysł na realizację niniejszego projektu zrodził się podczas wykonywania codziennych obowiązków zawodowych. W trakcie pracy nad jednym z projektów zaobserwowano interesującą tendencję związaną z procesem tworzenia testów modułowych. Znaczna część kodu testowego okazywała się redundantna, a wiele mechanizmów było wielokrotnie powielanych poprzez kopiowanie oraz wprowadzanie drobnych modyfikacji do istniejących fragmentów kodu.

Naturalnym rozwiązaniem tego problemu mogłoby wydawać się wydzielanie powtarzalnych elementów do funkcji pomocniczych lub bibliotek testowych oraz utrzymywanie ich w uporządkowanej i współdzielonej formie. Podejście to jednak wiąże się z istotnymi trudnościami, zarówno natury organizacyjnej, jak i technicznej. W praktyce często prowadzi ono do nadmiernego rozrostu bibliotek pomocniczych, niejednoznacznych interfejsów oraz problemów z ich długoterminowym utrzymaniem.

W szczególności pojawiają się następujące problemy:
\begin{enumerate}
  \item Jak określić moment, w którym dana funkcjonalność jest wystarczająco generyczna, aby mogła zostać wydzielona do wspólnej biblioteki?
  \item Jak zminimalizować ryzyko tworzenia funkcji lokalnie w plikach testowych, które z czasem zaczynają być wykorzystywane przez inne testy, prowadząc do niejawnych zależności pomiędzy teoretycznie niezależnymi modułami?
  \item Jak zapewnić spójny sposób tworzenia oraz utrzymania takich funkcji w projektach rozwijanych przez dziesiątki, a niekiedy setki programistów, o różnym poziomie doświadczenia i odmiennych nawykach pracy?
\end{enumerate}

Jednocześnie, jak powszechnie wiadomo, modele sztucznej inteligencji uczą się na podstawie zbiorów danych treningowych, a następnie wykorzystują wyuczone wzorce do rozwiązywania nowych problemów. Obserwacja ta stała się punktem wyjścia do rozważenia możliwości zastosowania modeli AI w obszarze automatycznego generowania wspólnych, powtarzalnych fragmentów kodu testowego na podstawie już istniejących przykładów.

\subsection{Założenia}

W odpowiedzi na postawione pytania badawcze rozważono zastosowanie modelu, który został wstępnie przeuczony na dostępnych danych, a następnie rozwijany wraz z pozyskiwaniem kolejnych zbiorów treningowych.

Na tym etapie pojawia się pierwszy istotny problem założeniowy, a mianowicie konieczność posiadania modelu sztucznej inteligencji. Dla uproszczenia przyjęto, że model SI dostarczany jest przez zewnętrzną firmę, natomiast zakres pracy obejmuje jego wykorzystanie, dalsze uczenie oraz wdrażanie. W tym kontekście pojawia się kluczowe pytanie niniejszej pracy: czy bardziej opłacalne jest oddelegowanie jednego pracownika do ręcznego utrzymywania i rozwijania tego oprogramowania, czy też lepszym rozwiązaniem jest automatyzacja tego procesu z wykorzystaniem narzędzi DevOps oraz pipeline’ów CI/CD?

W podejściu manualnym proces ponownego uczenia i wdrażania modelu realizowany jest przy znaczącym udziale człowieka. Obejmuje on ręczne przygotowanie danych, uruchamianie treningu, walidację wyników oraz publikację nowej wersji modelu. Taki sposób pracy, mimo swojej prostoty, wiąże się z ryzykiem błędów ludzkich, trudnościami w zachowaniu powtarzalności eksperymentów oraz ograniczoną skalowalnością procesu.

Alternatywą dla tego podejścia jest zastosowanie zautomatyzowanego pipeline’u CI/CD, w którym kolejne etapy cyklu życia modelu realizowane są w sposób powtarzalny i kontrolowany, z wykorzystaniem narzędzi znanych z obszaru DevOps. Automatyzacja umożliwia jednoznaczne powiązanie wersji modelu z wersją kodu źródłowego oraz danych treningowych, a także wprowadzenie obiektywnych kryteriów jakościowych decydujących o wdrożeniu nowej wersji modelu.

Wobec powyższych obserwacji zasadnym staje się przeanalizowanie opłacalności oraz efektywności automatyzacji procesu CI/CD w kontekście zarządzania cyklem życia modeli sztucznej inteligencji.

\subsubsection{Zakres projektu}

W ramach niniejszej pracy opracowano kompletny, w pełni zautomatyzowany system CI/CD dla modeli sztucznej inteligencji, oparty o rzeczywisty przypadek użycia: segmentację semantyczną obrazów wodomierzy z wykorzystaniem architektury U-Net. Wybór tego problemu jest celowy, ponieważ model segmentacji obrazów charakteryzuje się:

\begin{itemize}
  \item wyraźnie zdefiniowanymi metrykami jakości (współczynnik Dice, IoU),
  \item umiarkowanymi wymaganiami obliczeniowymi, umożliwiającymi trening w ramach dostępnego budżetu,
  \item dostępnością wstępnie wytrenowanego modelu bazowego, służącego jako punkt odniesienia dla bramki jakości,
\end{itemize}

Całość projektu została zaprojektowana w taki sposób, aby obejmować wszystkie kluczowe elementy wymagane do realizacji w pełni zautomatyzowanego procesu CI/CD dla modeli sztucznej inteligencji. Opracowany system uwzględnia konteneryzację komponentów z wykorzystaniem Docker, orkiestrację wdrożeń przy użyciu Kubernetes, a także izolację poszczególnych wersji modeli za pomocą namespace’ów. Proces obejmuje automatyczne wykrywanie zmian w kodzie źródłowym oraz danych treningowych, wersjonowanie modeli i danych (MLflow, DVC), a także kontrolę jakości modeli na podstawie zdefiniowanych metryk.

Integralną częścią rozwiązania jest środowisko chmurowe, w którym realizowane są procesy uczenia, testowania, wdrażania modeli SI oraz możliwości predykcji przez użytkownika we wdrożonym modelu. W ramach projektu wdrożono mechanizmy monitoringu infrastruktury oraz aplikacji (Prometheus, Grafana), umożliwiające analizę zużycia zasobów oraz stabilności działania systemu. Cały pipeline CI/CD został zaprojektowany w sposób powtarzalny i możliwy do odtworzenia, co pozwala na bezpośrednie porównanie z podejściem manualnym oraz ocenę wpływu automatyzacji na czas wdrożenia, niezawodność i koszty operacyjne.
